\documentclass[11pt]{article}
\usepackage[margin=1in]{geometry}
\usepackage{amsmath,amssymb,amsthm,enumitem,hyperref}
\hypersetup{colorlinks=true,linkcolor=blue,urlcolor=blue,citecolor=blue}
\newtheorem{theorem}{Theorem}
\newtheorem{lemma}{Lemma}
\newcommand{\Z}{\mathbb Z}
\newcommand{\saw}[1]{\left(\!\left(#1\right)\!\right)}
\title{Divisibility of Shibukawa's elliptic Dedekind sum\\\large MathDB 342271}
\author{Principia Math In Progress}
\date{September 18, 2026}
\begin{document}
\maketitle
\begin{abstract}
We prove Conjecture 13 of Shibukawa: for coprime even $a$ and positive
odd $b$, the integer $bQ(a;b)=4bs_i(a;b)$ is divisible by $4$, and by
$12$ when $3\nmid b$. The proof uses Shibukawa's established expression
in classical Dedekind sums and the classical Meyer inversion formula.
A square multiplier has positive permutation sign, which gives the
factor $4$. A separate elementary denominator argument gives the
factor $3$. We prove a more general three-term congruence for every
modulus, including even moduli.
\end{abstract}

\section*{Task 1. Full solution}
The domain inherited from \cite{S} is $a\in\Z\setminus\{0\}$,
$b\in\Z_{>0}$, $\gcd(a,b)=1$, and $a+b$ odd. We retain these
hypotheses in interpreting the question. Write $s_i(a;b)$ for the
elliptic sum in \cite{S} evaluated at $\tau=i=\sqrt{-1}$ and $Q(a;b)$
for its rational part. In the normalization of that paper,
\[
 s_\tau(a;b)=Q(a;b)\left(\frac13-\frac16\lambda(\tau)\right),
 \qquad \lambda(i)=\frac12,\qquad 4s_i(a;b)=Q(a;b).                 \tag{1}
\]
\begin{theorem}\label{target}
For every even integer $a\ne0$ and positive odd integer $b$ with
$\gcd(a,b)=1$,
\[
 \boxed{\ bQ(a;b)=4bs_i(a;b)\in4\Z.\ }
\]
If $3\nmid b$, then $bQ(a;b)\in12\Z$. In particular the latter
conclusion holds for every positive modulus $b=6m\pm1$.
\end{theorem}
This is the full divisibility conjecture, including its stronger
claim for the moduli $6m\pm1$, in both the preprint and the published
final of \cite{S}.

Define the sawtooth function by
\[
 \saw{x}=\begin{cases}x-\lfloor x\rfloor-\tfrac12,&x\notin\Z,\\
                              0,&x\in\Z,\end{cases}
 \qquad
 s(h;b)=\sum_{j=1}^{b-1}\saw{\frac jb}\saw{\frac{hj}b}.             \tag{2}
\]
Here $b\ge1$, $h$ is any integer, and an empty sum is zero. In
particular (2) gives the usual extension to noncoprime arguments.

\subsection*{Reduction to three classical Dedekind sums}
Shibukawa's formula \cite[(4.3)]{S}, obtained using
Sitaramachandra Rao's Hardy-sum identity \cite{R}, is
\[
 Q(a;b)=6\bigl(s(a;2b)+s(2a;b)-2s(a;b)\bigr).                    \tag{3}
\]
This identity is an established input. Put $a=2c$. Directly from (2),
\[
 s(2c;2b)=s(c;b).                                                \tag{4}
\]
Indeed, for $1\le j<b$ the terms indexed by $j$ and $j+b$ have the
same second sawtooth factor, whereas their first factors satisfy
\[
 \saw{\frac{j}{2b}}+\saw{\frac{j+b}{2b}}
       =\frac jb-\frac12=\saw{\frac jb}.
\]
The term indexed by $b$ vanishes. Thus (3) becomes
\[
 bQ(2c;b)=6b\bigl(s(c;b)+s(4c;b)-2s(2c;b)\bigr).                 \tag{5}
\]

\subsection*{A congruence for every unit multiplier}
\begin{theorem}\label{general}
Let $b\ge1$ and $h,t\in\Z$ satisfy $\gcd(ht,b)=1$. Then
\[
 D_b(h,t):=6b\bigl(s(h;b)+s(ht^2;b)-2s(ht;b)\bigr)\in4\Z.
                                                                    \tag{6}
\]
If $3\nmid b$, then $D_b(h,t)\in12\Z$.
\end{theorem}
We give a self-contained proof of the classical inversion identity
used in this theorem. It is Meyer's formula \cite{M}; see also
\cite[Theorem 1.2]{T} for its attribution.

\begin{lemma}[Meyer's inversion formula]\label{meyer}
Let $\gcd(h,b)=1$, and let $I(h,b)$ be the inversion number of
the permutation $\sigma_h:j\mapsto hj\pmod b$ of
$\{0,\ldots,b-1\}$, using least nonnegative residues. Then
\[
 I(h,b)=\frac{(b-1)(b-2)}4-3b\,s(h;b).                            \tag{7}
\]
\end{lemma}
\begin{proof}
The case $b=1$ is immediate. Suppose $b\ge2$ and let $r_j$ be the
least nonnegative residue of $hj$. The values $r_1,\ldots,r_{b-1}$
permute $1,\ldots,b-1$. For $0\le i<j\le b-1$,
\[
 r_{j-i}=r_j-r_i+b\,\boldsymbol{1}_{\{r_i>r_j\}}.
\]
Summing this equality over all such pairs gives
\begin{align*}
 bI(h,b)
  &=\sum_{k=1}^{b-1}(b-k)r_k
       -\sum_{k=1}^{b-1}(2k-b+1)r_k\\
  &=\frac{b(2b-1)(b-1)}2-3\sum_{k=1}^{b-1}k r_k.
\end{align*}
On the other hand, expanding (2) gives
\[
 s(h;b)=\frac1{b^2}\sum_{k=1}^{b-1}k r_k-\frac{b-1}{4}.           \tag{8}
\]
Eliminating the residue sum proves (7).
\end{proof}

\begin{proof}[Proof of Theorem~\ref{general}]
For $b=1$ all sums vanish, so assume $b\ge2$. Multiplication of
residue classes gives
\[
 \sigma_{ht^2}=\sigma_t\circ\sigma_t\circ\sigma_h.
\]
Since a square of a permutation has positive sign,
$(-1)^{I(ht^2,b)}=(-1)^{I(h,b)}$. Consequently
\[
 I(ht^2,b)\equiv I(h,b)\pmod2.                                   \tag{9}
\]
In (6), Meyer's identity cancels the constant term and yields
\[
 D_b(h,t)=4I(ht,b)-2I(h,b)-2I(ht^2,b).                             \tag{10}
\]
The last two inversion numbers have even sum by (9). Thus (10)
is an integer divisible by $4$.

For the additional factor $3$, suppose $3\nmid b$. For any
$u$ coprime to $b$, let $r_j$ now be the residue of $uj$. The identity
\[
 2\sum_{j=1}^{b-1}j^2=\frac{b(b-1)(2b-1)}3
\]
and $3\mid(b-1)(2b-1)$ imply $b\mid2\sum j^2$.
Since $r_j\equiv uj\pmod b$, also $b\mid2\sum j r_j$. By (8),
\[
 2b\,s(u;b)=\frac2b\sum_{j=1}^{b-1}j r_j-\frac{b(b-1)}2\in\Z.
                                                                    \tag{11}
\]
It follows that each of $6b\,s(h;b)$, $6b\,s(ht;b)$, and
$6b\,s(ht^2;b)$ belongs to $3\Z$, so $3\mid D_b(h,t)$.
Together with the factor $4$ already proved, this proves the theorem.
\end{proof}

\begin{proof}[Proof of Theorem~\ref{target}]
For $a=2c$ and odd $b$ with $\gcd(a,b)=1$, both $c$ and $2$ are
units modulo $b$. Take $h=c$ and $t=2$ in Theorem~\ref{general}
and use (5). Equation (1) identifies the elliptic formulation.
The proof uses residue classes and therefore also covers negative
$a$. For $b=1$, (3)--(4) give $Q(a;1)=0$.
\end{proof}

The uniform factors cannot be increased: exact values are
\[
 3Q(2;3)=4,\qquad 5Q(4;5)=12,\qquad 7Q(2;7)=12.
\]
Thus the factor $4$ is sharp over all odd moduli, and the factor
$12$ is sharp even within each of the two classes $6m-1$ and $6m+1$.
The contribution here is the application of permutation-sign
multiplicativity to (3), with the denominator calculation (11).
Formula (3) and Meyer's identity are explicitly credited prior work.

\section*{Tasks 2--5. Proposed follow-ups}
\begin{enumerate}[label=\textbf{Task \arabic*.},start=2,leftmargin=*]
\item For nonzero $bQ(a;b)$, determine its exact $2$-adic valuation
from the Euclidean algorithm for $(a/2,b)$, and obtain a useful
criterion for divisibility by $8$ or higher powers of $2$.
\item When $3\mid b$, determine the precise condition for
$12\mid bQ(a;b)$ and the higher $3$-adic valuations of its nonzero
values.
\item For each fixed odd $b>1$, determine the greatest common
divisor of the values $bQ(2c;b)$ over all units $c$ modulo $b$,
including any exceptional moduli at which all these values vanish.
\item Study denominator and divisibility bounds for the higher
elliptic sums introduced by Shibukawa. Determine whether analogous
permutation statistics control the rational coefficients in these
Dedekind--Apostol sums.
\end{enumerate}

\paragraph{Verification.}
Four exact arithmetic groups pass. An independent Hardy-sum formula
checks 8,151 eligible pairs; 1,736 permutations and 61,854 unit pairs
check inversions and signs directly; 140,390 triples check the general
theorem, supplemented by 26 larger cases; and 3,157 signed pairs check
negative arguments and period, with the boundary and sharp values.
These finite checks support the proof and do not replace it.
Review is by the authoring agent only.

\begin{thebibliography}{9}
\bibitem{S} G. Shibukawa,
\emph{Some arithmetic properties of an elliptic Dedekind sum},
JP J. Algebra Number Theory Appl. \textbf{46}(1) (2020), 35--54,
\href{https://doi.org/10.17654/NT046010035}{doi:10.17654/NT046010035};
\href{https://arxiv.org/abs/1903.12200}{arXiv:1903.12200v2} (2019).
Conjecture 13 is on page 52 of the final, page 13 of the preprint.
\bibitem{R} R. Sitaramachandra Rao,
\emph{Dedekind and Hardy sums}, Acta Arith. \textbf{48}(4) (1987),
325--340. Identity quoted as Lemma 7 in \cite{S}.
\bibitem{M} C. Meyer,
\emph{\"Uber einige Anwendungen Dedekindscher Summen},
J. Reine Angew. Math. \textbf{198} (1957), 143--203.
\bibitem{T} E. Tsukerman,
\emph{Equality of Dedekind sums mod $\Z,2\Z$ and $4\Z$},
\href{https://arxiv.org/abs/1408.3225}{arXiv:1408.3225v2} (2014),
Theorem 1.2.
\end{thebibliography}
\end{document}
